\section{Circles tangent to consecutive sides}\label{sec:side-tangent}

There is a second discrete curvature attached to a convex polygon: the
reciprocal radius of the inward circle tangent to three consecutive
supporting lines. Wegner's four-edge theorem concerns these radii
\cite[Theorem~1]{Wegner1994}; see also \cite[Corollary~21.11]{Pak2010}. We give a converse as a companion to
Theorem~\ref{thm:main}. The construction below uses a tangential polygon and
positive edge scaling; it is independent of the four-length system used for
the main theorem.

\begin{definition}\label{def:side-tangent-circle}
Let $P_i$ be the vertices of a counterclockwise strictly convex polygon,
with outward unit normal $N_i$ on edge $[P_i,P_{i+1}]$.
The \emph{inward three-line circle at edge $i$} has center $C_i$ and radius
$r_i>0$ satisfying
\begin{equation}\label{eq:side-tangent-distances}
 \langle N_{i-1},P_i-C_i\rangle
 =\langle N_i,P_i-C_i\rangle
 =\langle N_{i+1},P_{i+1}-C_i\rangle=r_i.
\end{equation}
Its central contact is required to lie in the open edge $(P_i,P_{i+1})$.
The contacts on the two neighbouring supporting lines may lie on their
extensions. The circle need not be contained in the entire polygon.
\end{definition}
% Lean: Gluck.Discrete.SideTangent.ThreeLineCircle
% Lean: Gluck.Discrete.SideTangent.IsSideTangentRealization

\begin{figure}[htbp]
\centering
\figfile{side-tangent-circle}{0.52\textwidth}
\caption{The selected circle is tangent to three supporting lines, with
its central contact in the open edge. Here one neighbouring contact lies
beyond its edge. All coordinates, circle equations and contact parameters
in this example are checked with exact rational arithmetic.}
\label{fig:side-tangent-circle}
\end{figure}

Put $c_i=r_{i-1}-r_i$. By four cyclic sign changes we mean that there are
$i_1<i_2<i_3<i_4<i_1+n$ with
$c_{i_1}>0$, $c_{i_2}<0$, $c_{i_3}>0$ and $c_{i_4}<0$; zero differences are
retained but do not count as a sign. This formulation includes repeated
radii without a genericity assumption.

\begin{theorem}\label{thm:side-tangent-converse}
Let $n\ge4$ and let $r_i>0$ be a cyclic sequence. Suppose that $r$ is
constant or that $c_i=r_{i-1}-r_i$ has four cyclic sign changes. Then a
simple strictly convex polygon realizes $r_i$ as the radii of its inward
three-line circles. Each center and radius are uniquely determined by the
three equations~\eqref{eq:side-tangent-distances}, and every central contact
lies in its open edge.
\end{theorem}
% Lean: Gluck.Discrete.SideTangent.sideTangentConverse
% Lean: Gluck.Discrete.SideTangent.IsSideTangentRealization.tangent_sphere

The main ingredient is a balanced polygon circumscribed about the unit
circle. Its vertices, rather than the contact points, carry the coefficients.

\begin{lemma}\label{lem:side-tangent-balance}
Suppose $\sum_i c_i=0$ and $c$ has four cyclic sign changes. There is a
strictly convex polygon $Y_i$, circumscribed about the unit circle with each
contact in its open edge, such that $\sum_i c_iY_i=0$.
\end{lemma}
% Lean: Gluck.Discrete.SideTangent.exists_balanced_cyclicFan_of_dfv

\begin{proof}
Choose four alternating nonzero coefficients, and place their labels at
$(1,1),(-1,1),(-1,-1),(1,-1)$ in that order. Collapse all intervening labels
to the corresponding side contact of the square. All labels and their
coefficients remain present. Denote this limiting configuration by $X_i$,
and consider the finite rational field
\[
 G(a)=\sum_i\frac{c_i X_i}{1-\langle a,X_i\rangle}.
\]
For sufficiently small $\varepsilon>0$, all denominators are positive on
$D_\varepsilon=\{(x,y):|x|+|y|\le1-\varepsilon\}$. On its upper boundary arc
$y\ge |x|$, one of the two upper corner denominators equals $\varepsilon$.
If the alternating corner coefficients are $\alpha,-\beta,\gamma,-\delta$
with $\alpha,\beta,\gamma,\delta>0$, and the left and right contact
coefficients are $u,v$, then
\[
 \mathop{\rm Re}G(x,y)\ge
 \frac{\min(\alpha,\beta)}{\varepsilon}
       -\gamma-\delta-2|u|-2|v|>0.
\]
The upper and lower contact labels have zero horizontal contribution,
including near their poles. The corresponding bounds on the other three
arcs give, successively, $\mathop{\rm Im}G<0$ on the left,
$\mathop{\rm Re}G<0$ on the bottom and $\mathop{\rm Im}G>0$ on the right.
Thus the boundary field has winding number $-1$.

Now separate the repeated normals into small, strictly ordered clusters.
This produces a genuine tangential polygon $X_i^{\eta}$, with every edge
positive and every contact interior, converging to the labelled square as
$\eta\downarrow0$. Positivity of all denominators and the strict boundary
margins persist uniformly on the compact diamond. The perturbed field has
the same nonzero winding number and hence a zero $a$ in $D_\varepsilon$.
In particular $|a|<1$.

Set $\sigma=\sqrt{1-|a|^2}$ and
\[
 A=\sigma I+\frac{aa^{\mathsf T}}{1+\sigma},\qquad
 B_a(x)=\frac{Ax-a}{1-\langle a,x\rangle}.
\]
This projective transformation preserves the unit disk, its tangent lines,
and the cyclic order and open contacts of our polygon. All relevant
denominators are positive. Since $\sum_i c_i=0$,
\[
 \sum_i\frac{c_i}{1-\langle a,X_i^{\eta}\rangle}
 =\langle a,G(a)\rangle=0,
 \qquad
 \sum_i c_i B_a(X_i^{\eta})=A G(a)-a\langle a,G(a)\rangle=0.
\]
The vertices $Y_i=B_a(X_i^{\eta})$ have the required properties.
\end{proof}

\begin{lemma}\label{lem:side-tangent-scaling}
Let $Y_i$ be a counterclockwise strictly convex polygon circumscribed about the unit circle,
with its contacts in its open edges. If $r_i>0$ and
$\sum_i(r_{i-1}-r_i)Y_i=0$, the cyclic edge equations
\begin{equation}\label{eq:side-tangent-scaling}
 P_{i+1}-P_i=r_i(Y_{i+1}-Y_i)
\end{equation}
define a closed simple strictly convex polygon. Its inward three-line
circle at edge $i$ has center $C_i=P_i-r_iY_i$ and radius $r_i$.
\end{lemma}
% Lean: Gluck.Discrete.SideTangent.sum_scaled_edges
% Lean: Gluck.Discrete.SideTangent.strictSupport_of_positive_edge_scaling
% Lean: Gluck.Discrete.SideTangent.realizesSideTangent_of_scaledFan

\begin{proof}
Cyclic summation by parts gives
$\sum_i r_i(Y_{i+1}-Y_i)=\sum_i(r_{i-1}-r_i)Y_i=0$, so the edges close.
Positive scaling preserves their ordered directions and their total turn
of $2\pi$. For any fixed edge, the projections of the subsequent edges on
its left normal are first positive and then negative. Every intermediate
partial sum is positive, using closure for the second part. Hence every
nonincident vertex is strictly to its left. This gives strict support and
simplicity.

The outward normals satisfy
$\langle N_{i-1},Y_i\rangle=\langle N_i,Y_i\rangle=1$ and
$\langle N_{i+1},Y_{i+1}\rangle=1$. Substituting
$C_i=P_i-r_iY_i=P_{i+1}-r_iY_{i+1}$ gives all three
signed-distance equations. The central contact is the image of $N_i$ under
the positive homothety $x\mapsto C_i+r_ix$, and therefore lies in its open
edge. Adjacent normals are linearly independent. The first two equations
force any other center with radius $s$ to equal $P_i-sY_i$; the third then
gives $s=r_i$, since $\langle N_{i+1},Y_i\rangle<1$. Finally, unit normals
turn the signed-distance equations into actual circle tangency.
\end{proof}

\begin{figure}[tb]
\centering
\figfile{side-tangent-construction}{\textwidth}
\caption{A balanced polygon circumscribed about the unit circle (left) and
its positive edge scaling (right). The prescribed cyclic radii are
$(1,2,1,3/2,1,2,1,3/2)$. Closure, strict support and all eight circles are
verified exactly. The highlighted upper edge is the central edge of the
circle in Figure~\ref{fig:side-tangent-circle}.}
\label{fig:side-tangent-construction}
\end{figure}

\begin{proof}[Proof of Theorem~\ref{thm:side-tangent-converse}]
For nonconstant data, apply Lemma~\ref{lem:side-tangent-balance} to
$c_i=r_{i-1}-r_i$, whose sum is zero, and then
Lemma~\ref{lem:side-tangent-scaling}. For constant $r_i=r$, take a regular
polygon circumscribed about the unit circle and scale its edges by $r$.
\end{proof}

The forward implication is Wegner's four-edge theorem
\cite[Theorem~1]{Wegner1994}, whose convention already includes plateau
extrema. For the supporting-line circles of
Definition~\ref{def:side-tangent-circle}, we give a proof using the signed
affine-balance argument of \cite[Theorem~21.15]{Pak2010}, allowing zero
coefficients.

\begin{lemma}\label{lem:side-tangent-signs}
Let $Y_i$ be the cyclic vertices of a counterclockwise strictly convex polygon. If
$\sum_i c_i=0$ and $\sum_i c_iY_i=0$, then either every $c_i=0$ or the cyclic
coefficient sequence has at least four sign changes.
\end{lemma}
% Lean: Gluck.Discrete.SideTangent.zero_or_dfv_of_strict_affine_balance

\begin{proof}
A nonzero zero-sum sequence has both signs. If it has no four alternating
signs, it has one nonpositive block and one nonnegative block, with zeros
allowed in either. Choose a negative entry $c_a$ beginning the first block
and a positive entry $c_b$ beginning the second, using cyclic lifts
$a<b<a+n$. The chord $Y_aY_b$ has strictly negative signed area against
vertices in $(a,b)$ and strictly positive signed area against those in
$(b,a+n)$. Thus
\[
 c_i[Y_a,Y_b,Y_i]\ge0\quad\hbox{for every }i,
 \qquad \sum_i c_i[Y_a,Y_b,Y_i]=0,
\]
where $[A,B,C]$ is the signed determinant. The equality follows from the
zero mass and zero vector sum, since the determinant is affine in its third
argument. Every summand vanishes. Strict chord separation forces all
coefficients except $c_a,c_b$ to vanish. The two balance equations then give
$c_a(Y_a-Y_b)=0$, contradicting $c_a<0$ and distinctness of the vertices.
\end{proof}

\begin{corollary}\label{cor:side-tangent-iff}
Let $n\ge4$ and $r_i>0$. The sequence $r$ is realized by the inward
three-line circles of a simple strictly convex polygon, with their central
contacts in the open edges, if and only if $r$ is constant or its cyclic
differences $r_{i-1}-r_i$ have at least four sign changes.
\end{corollary}
% Lean: Gluck.Discrete.SideTangent.sideTangentForward
% Lean: Gluck.Discrete.SideTangent.realizesSideTangent_iff_admissible

\begin{proof}
Sufficiency is Theorem~\ref{thm:side-tangent-converse}. For necessity, start
from the actual polygon and its circles, and set
$Y_i=(P_i-C_i)/r_i$. The two incident normal equations give
$\langle N_{i-1},Y_i\rangle=\langle N_i,Y_i\rangle=1$.
These normals are linearly independent, because the polygon is strictly
convex. The circle at the preceding edge satisfies the same two normalized
equations at $P_i$, so it gives the same $Y_i$. Consequently
\[
 Y_{i+1}-Y_i=\frac{P_{i+1}-P_i}{r_i}.
\]
This is a closed positive edge scaling of the original polygon; the strict
support argument in Lemma~\ref{lem:side-tangent-scaling} therefore applies to
$Y$. Summation by parts gives
$\sum_i(r_{i-1}-r_i)Y_i=0$, and the differences themselves sum to zero.
Lemma~\ref{lem:side-tangent-signs} supplies four sign changes unless every
difference vanishes. In the latter case the cyclic radius sequence is constant.
\end{proof}
